\chapter{Discussion}

\section{Summary of Main Findings}
This survey of 231 female university students in Dhaka found that four in five (80.3\%) accepted HPV vaccination, either because they had already received a dose (17.0\%) or because they were likely or very likely to be vaccinated. Yet knowledge was low (mean 6.32 of 15; two thirds with poor knowledge), uptake was small, and the commonest reason for not being vaccinated was simply not knowing about the vaccine. Having heard of HPV and holding a more positive attitude were the factors independently associated with acceptance. Healthcare and non-healthcare students differed in knowledge and attitude scores but not in acceptance.

\section{Level of Acceptance}
Acceptance was high. Among students who were unvaccinated or unsure, 76.3\% said they were likely or very likely to be vaccinated, which is very close to the 76.6\% of unvaccinated rural women~\cite{banik2022} and the 77\% of unvaccinated educated women~\cite{taniya2023} who said they were willing in earlier Bangladeshi studies. Willingness thus appears similar in quite different groups of women. Only about one in five unvaccinated students (22.1\%) were uncertain and almost none refused (1.6\%), so rejection of the vaccine was not the main obstacle.

Actual uptake in the sample (39 students, 17.0\% with at least one dose) was higher than the 5.3\% to 16.1\% reported in other Bangladeshi studies of women and students~\cite{banik2022,taniya2023,chowdhury2022,promi2025}, and higher than the 9.5\% among female students in Beijing~\cite{liu2020}. This should be read with caution. Because the 2023--2024 national campaign was aimed at girls aged 10--14 years, no respondent was in the target group when it was run, so any vaccination in this sample was most likely obtained privately or outside the programme. A further 19 students (8.2\%) were not sure whether they had been vaccinated, which shows how little contact many students have with their own vaccination records.

\section{Awareness, Information Sources and Knowledge}
Most students (79.8\%) had heard of HPV, similar to the 72.6\% in Beijing~\cite{liu2020}, and awareness was higher among healthcare students (90.5\% against 72.2\%). Awareness alone, however, did not mean understanding. The mean knowledge score was 42\% of the maximum. Students knew what HPV stands for and that it can lead to cervical cancer, but most were unaware that HPV often needs no treatment, that the vaccine is most effective before sexual activity begins, and that males can also benefit. Few correctly rejected the claim that the vaccine protects against all strains, and many answered ``not sure''. A similar uncertainty about the range of protection was reported among students at Jahangirnagar University~\cite{emon2025}. The pattern points to shallow rather than absent knowledge: students have encountered the terms but have not learned the details that matter for a vaccination decision.

The curriculum was the commonest first source of information (54.9\%), followed by social media (30.8\%). Health campaigns (7.7\%) and healthcare providers (9.3\%) were rarely named. This fits the finding that over half of the unvaccinated students did not know about the vaccine, and it suggests that campaign messages have not reached young adults. It also fits the preference reported by students in Beijing for college education as the main source of vaccine information~\cite{liu2020}. Because social media is the second most common source, accuracy of the material shared there also matters.

\section{Healthcare and Non-Healthcare Students}
Healthcare students had better knowledge (7.42 against 5.54) and slightly more favourable attitude scores, but acceptance was practically the same in both groups (79.8\% against 80.7\%). Even among healthcare students, only 11.6\% reached the good knowledge category and 61.1\% had poor knowledge. Medical and health students in other Bangladeshi studies also showed gaps between knowledge and practice~\cite{chowdhury2022,promi2025}. The similarity in acceptance suggests that acceptance is driven less by academic training than by awareness of the vaccine and by general attitude. It also suggests that HPV and vaccine content in healthcare curricula may need to be strengthened.

\section{Attitudes}
Attitudes were overwhelmingly positive: 88.3\% had a positive attitude, and more than 90\% agreed that young women should be vaccinated. The gap between positive attitudes and low knowledge, which might be called a knowledge--attitude gap, is consistent with the finding that attitude was the strongest predictor of willingness in a large student survey~\cite{emon2025}. Attitudes about screening were more lukewarm: only 64.8\% agreed that Pap smear testing is necessary and a third were neutral, so preventive information on screening deserves attention alongside vaccination. Support for educational measures was very high, with 90.3\% in favour of including HPV vaccination in institutional health programmes and 83.1\% in favour of compulsory HPV education.

\section{Factors Associated With Acceptance}
Having heard of HPV was the strongest and most stable correlate of acceptance, with about three times the odds in both adjusted models and in the sensitivity analysis. Students who had never heard of HPV were far less likely to accept (63.0\% against 85.0\%). This is plausible: a vaccine cannot be wanted by someone who does not know it exists. Attitude score was also independently associated with acceptance, but its association is partly built in, since several attitude items (for example, ``all young women should be vaccinated'') are close to the outcome itself. The knowledge score was associated with acceptance in unadjusted analysis but lost significance once awareness was included, which suggests that awareness captures most of its effect. Age, academic stream, parents' education, prior STI education and parental discussion were not significantly associated with acceptance.

Cost and access recurred across the results.
Students rated lack of awareness, vaccine cost, lack of provider recommendation and limited availability as the main barriers.
Nearly half (45.9\%) would be encouraged by a free or subsidised vaccine and 42.4\% by a doctor's recommendation.
A further 82.0\% of respondents said that provider advice would increase their likelihood of vaccination. Provider recommendation is also a predictor of uptake among students in China~\cite{liu2020}. Cultural or religious concerns, by contrast, were rated lowest as a barrier, and only one student gave family or religious opposition as a reason for not being vaccinated, even though about 69\% said that discussing HPV is taboo in their social context. Taboo may reduce conversation and information, rather than lead to outright rejection of the vaccine.

\section{Implications}
\begin{itemize}[leftmargin=*,itemsep=2pt]
\item \textbf{Universities.} Because most students already support vaccination, campus awareness activities should concentrate on telling students that the vaccine exists, who can take it and what it protects against, and on correcting specific misconceptions.
\item \textbf{Curricula.} HPV, cervical cancer and screening content should be taught to students in all faculties, and clinical vaccine counselling should be strengthened for healthcare students.
\item \textbf{Health providers and authorities.} Doctors should routinely recommend the vaccine to young women. Because students named cost and availability, catch-up vaccination at low or no cost, ideally at campus health centres, deserves consideration.
\item \textbf{Communication.} Social media is a major information source, so accurate material from health professionals should be placed there.
\end{itemize}

\section{Strengths and Limitations}
\subsection*{Strengths}
The study measured acceptance directly, combining vaccination status with future intention, and examined it alongside knowledge, attitude, sources of information, barriers and motivators. The knowledge and attitude scales had good internal consistency (KR-20 = 0.828; $\alpha$ = 0.823), the sample included students from several academic fields, and the main finding held in a sensitivity analysis.

\subsection*{Limitations}
\begin{enumerate}[leftmargin=*,itemsep=2pt]
\item The design was cross-sectional, so associations cannot be taken as cause and effect, and stated intention may not translate into vaccination.
% TODO: Confirm the final sampling approach and the institutions represented.
\item The sample was a convenience sample from a small number of Dhaka universities \confirm{sampling approach}, so the findings may not apply to students elsewhere in Bangladesh, or to those in rural areas or other kinds of institution. The institution was not recorded for 87 respondents.
\item Only 45 students were classed as not accepting. The regression estimates are therefore imprecise, with wide confidence intervals, and some real associations may not have been detected. Model fit was modest.
\item The outcome combined past vaccination with future intention. Intention was reported in hypothetical terms, and 19 students did not know whether they had been vaccinated.
\item All data were self-reported and may be affected by social desirability, which could have inflated positive attitudes. ``Not sure'' answers were scored as incorrect in the knowledge score, which can lower scores compared with studies that score them differently.
\item The attitude score contains items that overlap conceptually with acceptance, so its association with the outcome should be interpreted with care.
\item Many comparisons were made without correction for multiple testing.
\end{enumerate}
